\chapter{Tacrolimus Level}

\section*{What Is This Test?}
The tacrolimus level (also known as FK506 level) measures the whole blood concentration of tacrolimus, a calcineurin inhibitor immunosuppressant that is the most widely used maintenance immunosuppressant in solid organ transplantation today. Tacrolimus is 10--100 times more potent than cyclosporine on a per-weight basis. It binds to FK-binding protein 12 (FKBP12), and the tacrolimus-FKBP12 complex inhibits calcineurin, preventing NFAT dephosphorylation and nuclear translocation, thereby blocking IL-2 gene transcription and T-cell activation. Tacrolimus is metabolized primarily by CYP3A4/CYP3A5 in the liver and intestine, and it is a substrate of P-glycoprotein (ABCB1). It has a narrow therapeutic index. Subtherapeutic levels risk acute rejection; supratherapeutic levels cause nephrotoxicity (similar to cyclosporine), neurotoxicity (tremor --- the most common symptom, headache, insomnia, paresthesias, seizures, PRES), hyperglycemia and new-onset diabetes after transplantation (NODAT, more common than with cyclosporine), hyperkalemia, hypomagnesemia, hypertension, and alopecia. Unlike cyclosporine, tacrolimus does NOT cause gingival hyperplasia or hirsutism. TDM uses whole blood (EDTA) collected as a 12-hour trough (C0) level. Tacrolimus is extensively distributed in erythrocytes (85--95\% of whole blood concentration), requiring whole blood measurement; serum or plasma levels are inaccurate.

\section*{Alternative Names}
FK506 Level; Tacrolimus Blood Level; Prograf Level; Advagraf Level; Astagraf XL Level; Protopic Level; Tacrolimus TDM

\section*{Principle}
Tacrolimus is measured by LC-MS/MS (gold standard, no metabolite cross-reactivity) or by automated chemiluminescent microparticle immunoassay (CMIA, Abbott Architect/Abbott Alinity), or enzyme multiplied immunoassay technique (EMIT). Immunoassays cross-react with inactive tacrolimus metabolites (13-O-demethyl, 15-O-demethyl, 31-O-demethyl tacrolimus) and overestimate the tacrolimus concentration by 10--25\% compared to LC-MS/MS, most significantly in patients with hepatic dysfunction who accumulate metabolites. Specimen: whole blood in EDTA (lavender-top). Level is a 12-hour trough (immediately before morning dose).

\section*{History}
Tacrolimus was discovered in 1984 from the fermentation broth of the soil bacterium Streptomyces tsukubaensis, found in a soil sample collected near Tsukuba, Japan. The discovery was made by the Fujisawa Pharmaceutical Company (now Astellas Pharma). It was first used clinically in liver transplantation by Thomas Starzl in 1989. FDA approval was granted in 1994 for liver transplantation and 1997 for kidney transplantation. Tacrolimus has largely replaced cyclosporine as the preferred calcineurin inhibitor for most transplant recipients due to lower acute rejection rates and a more favorable side effect profile (less hirsutism and gingival hyperplasia, though more neurotoxicity and NODAT).

\section*{Why This Test Is Ordered}
Routine TDM in transplant recipients to maintain levels within target ranges. Trough levels are typically: 8--12 ng/mL (first 3 months post-transplant, higher immunosuppression), 5--8 ng/mL (maintenance >3 months), and as low as 3--5 ng/mL (late maintenance with no rejection history). Frequent monitoring when initiating or discontinuing CYP3A4-interacting drugs. Investigation of suspected nephrotoxicity or neurotoxicity. Suspected non-adherence.

\section*{Primary vs Secondary Test}
TDM is a secondary monitoring test. Tacrolimus dosing is primarily managed by clinical response, graft function, and tolerability; the whole blood trough level is the essential secondary tool used to individualize the dose within transplant-specific target ranges and to help distinguish rejection from calcineurin-inhibitor nephrotoxicity.

\section*{How This Test Is Performed}
A blood sample is collected by standard venipuncture into an EDTA (lavender-top) tube; whole blood is required because tacrolimus is concentrated in erythrocytes, and serum or plasma levels are inaccurate. Tacrolimus is measured by LC-MS/MS (the gold standard, with no metabolite cross-reactivity) or by automated immunoassay (CMIA or EMIT), which cross-reacts with inactive metabolites and typically overestimates concentrations by 10--25\% compared with LC-MS/MS. Results are available within a few hours for urgent requests and 1--2 days for routine samples; the level is reported as a 12-hour trough.

\section*{What Preparation Is Needed}
The sample should be drawn as a 12-hour trough, immediately before the morning dose, at steady state (3--5 days after initiation or dose change), with dosing at consistent 12-hour intervals. Because tacrolimus is a CYP3A4/P-gp substrate, the time of the last dose and any interacting drugs (azole antifungals, macrolides, calcium channel blockers raise levels; rifampin, phenytoin, carbamazepine lower levels) must be documented for correct interpretation.

\section*{Contraindications and Precautions}
\begin{itemize}
  \item No absolute contraindications to standard venipuncture
  \item Whole blood EDTA is mandatory; serum, plasma, and heparinized samples are not acceptable and produce inaccurate results
  \item Interpret the level in the context of time post-transplant, transplant organ, graft function, and concurrent interacting drugs; target ranges differ by protocol and center
  \item Immunoassay results may overestimate true concentrations in hepatic dysfunction (accumulated inactive metabolites); LC-MS/MS is preferred in this setting
  \item Nephrotoxicity, neurotoxicity, hyperkalemia, hypomagnesemia, and hyperglycemia can occur even within the target range, particularly early post-transplant, so the level alone does not exclude toxicity
\end{itemize}

\section*{Risks}
Minimal risk. Slight pain or bruising at the venipuncture site. Rarely: hematoma, infection, or vasovagal syncope.

\section*{What Happens After the Test}
The result is interpreted against the transplant-specific trough target (typically 8--12 ng/mL early post-transplant, 5--8 ng/mL maintenance, 3--5 ng/mL late maintenance). Subtherapeutic levels prompt a dose increase with repeat sampling at a new steady state; supratherapeutic levels with renal dysfunction prompt dose reduction and rechecking of the level and creatinine. Levels are typically repeated 2--3 days after any dose change and more frequently when interacting drugs are started or stopped.

\section*{Understanding Results}
Subtherapeutic: Risk of acute and chronic allograft rejection. Supratherapeutic: Nephrotoxicity (afferent arteriolar vasoconstriction, rising creatinine), neurotoxicity (tremor, headache, PRES), hyperglycemia/NODAT (tacrolimus directly inhibits insulin secretion from pancreatic beta-cells), hyperkalemia, hypomagnesemia (renal magnesium wasting, can exacerbate neurotoxicity and cardiac arrhythmias).

\section*{Normal Results}
Trough (C0, 12 hours post-dose, whole blood EDTA): 8--12 ng/mL (first 1--3 months), 5--8 ng/mL (maintenance), 3--5 ng/mL (late maintenance >1 year, rejection-free). Liver transplant: may target higher troughs (8--12 ng/mL). Tacrolimus levels should be interpreted with consideration of concomitant CYP3A4/P-gp modulating drugs, renal and hepatic function, and time post-transplant.

\section*{Follow-up Tests}
Serum creatinine/eGFR, serum electrolytes (potassium, magnesium, glucose), blood glucose/HbA1c (NODAT screening). Allograft biopsy if rejection is suspected or if renal dysfunction persists despite tacrolimus dose reduction. Drug interaction review (azole antifungals, macrolides, calcium channel blockers increase levels; rifampin, phenytoin, carbamazepine decrease levels).

\section*{References}
\begin{enumerate}
  \item Starzl TE, Todo S, Fung J, et al. FK 506 for liver, kidney, and pancreas transplantation. Lancet. 1989;2(8670):1000--1004.
  \item Ekberg H, Tedesco-Silva H, Demirbas A, et al. Reduced exposure to calcineurin inhibitors in renal transplantation (ELITE-Symphony). N Engl J Med. 2007;357(25):2562--2575.
  \item Staatz CE, Tett SE. Clinical pharmacokinetics and pharmacodynamics of tacrolimus in solid organ transplantation. Clin Pharmacokinet. 2004;43(10):623--653.
  \item Holt DW, Armstrong VW, Griesmacher A, Morris RG, Napoli KL, Shaw LM. International Federation of Clinical Chemistry/International Association of Therapeutic Drug Monitoring and Clinical Toxicology working group on immunosuppressive drug monitoring. Ther Drug Monit. 2002;24(1):59--67.
  \item Venkataramanan R, Swaminathan A, Prasad T, et al. Clinical pharmacokinetics of tacrolimus. Clin Pharmacokinet. 1995;29(6):404--430.
  \item Burtis CA, Ashwood ER, Bruns DE. Tietz Textbook of Clinical Chemistry and Molecular Diagnostics. 5th ed. St. Louis, MO: Elsevier Saunders; 2012.
\end{enumerate}

\clearpage
